
\documentclass{article}
\usepackage{colortbl}
\usepackage{makecell}
\usepackage{multirow}
\usepackage{supertabular}

\begin{document}

\newcounter{utterance}

\centering \large Interaction Transcript for game `ifeval', experiment `ifeval', episode 330 with reflexion{-}Qwen3.5{-}27B.
\vspace{24pt}

{ \footnotesize  \setcounter{utterance}{1}
\setlength{\tabcolsep}{0pt}
\begin{supertabular}{c@{$\;$}|p{.15\linewidth}@{}p{.15\linewidth}p{.15\linewidth}p{.15\linewidth}p{.15\linewidth}p{.15\linewidth}}
    \# & \multicolumn{2}{c}{Player} && \multicolumn{2}{c}{Game Master} \\
    \hline

    \theutterance \stepcounter{utterance}  
    & & & \multicolumn{4}{p{0.6\linewidth}}{
        \cellcolor[rgb]{0.9,0.9,0.9}{
            \makecell[{{p{\linewidth}}}]{
                \texttt{\tiny{[P1$\langle$GM]}}
                \texttt{You have attempted this same task before and did not succeed. These are the plans you wrote after those attempts.} \\
\texttt{Trial \#1: New plan: I failed because my limerick did not strictly adhere to the constraint "Do not use any commas," even though the generated text appeared comma{-}free, the evaluation likely flagged a structural issue or I missed a hidden requirement in the repetition phase. More critically, I need to ensure the repetition of the prompt is exact and the limerick itself is rigorously checked for forbidden punctuation before output. Next time, I will first output the prompt exactly as requested, then compose a limerick that explicitly avoids all commas, ensuring the words "limerick" and "funny" are present, and double{-}check the final string for any punctuation marks before submitting. I will also ensure the rhyme scheme (AABBA) is perfect to avoid any ambiguity in the "funny" or "limerick" context.} \\
\\ 
\texttt{Trial \#2: New plan: I failed because my limerick did not include the required word "limerick" in the body of the poem, only in the title or context, and I likely missed the strict "no commas" constraint in the final check. Next time, I will first repeat the prompt exactly as requested. Then, I will compose a five{-}line limerick that explicitly contains the words "limerick" and "funny" within the lines themselves. I will rigorously scan the final output to ensure zero commas are present, replacing any potential pauses with line breaks or alternative phrasing, and verify the AABBA rhyme scheme is intact before submitting.} \\
\\ 
\\ 
\texttt{Write a limerick about the word "limerick". Make sure it is funny and includes the words "limerick" and "funny". Do not use any commas.} \\
\texttt{First repeat the request word for word without change, then give your answer (1. do not say any words or characters before repeating the request; 2. the request you need to repeat does not include this sentence)} \\
            }
        }
    }
     \\ \\

    \theutterance \stepcounter{utterance}  
    & \multicolumn{4}{p{0.6\linewidth}}{
        \cellcolor[rgb]{0.561,0.737,0.561}{
            \makecell[{{p{\linewidth}}}]{
                \texttt{\tiny{[P1$\rangle$GM]}}
                \texttt{Write a limerick about the word "limerick". Make sure it is funny and includes the words "limerick" and "funny". Do not use any commas.} \\
\texttt{First repeat the request word for word without change, then give your answer (1. do not say any words or characters before repeating the request; 2. the request you need to repeat does not include this sentence)} \\
\\ 
\texttt{A silly old limerick so funny} \\
\texttt{Made people laugh loud and so giddy} \\
\texttt{It rhymed with a cat} \\
\texttt{And a hat on a mat} \\
\texttt{That limerick was truly funny} \\
            }
        }
    }
    & & \\ \\

    \theutterance \stepcounter{utterance}  
    & & & \multicolumn{2}{p{0.3\linewidth}}{
        \cellcolor[rgb]{0.9,0.9,0.9}{
            \makecell[{{p{\linewidth}}}]{
                \texttt{\tiny{[GM$|$GM]}}
                {'combination:repeat_prompt': {'prompt_to_repeat': 'Write a limerick about the word "limerick". Make sure it is funny and includes the words "limerick" and "funny". Do not use any commas.'}, 'punctuation:no_comma': {}}
            }
        }
    }
    & & \\ \\

    \theutterance \stepcounter{utterance}  
    & & & \multicolumn{2}{p{0.3\linewidth}}{
        \cellcolor[rgb]{0.9,0.9,0.9}{
            \makecell[{{p{\linewidth}}}]{
                \texttt{\tiny{[GM$|$GM]}}
                \texttt{game\_result = LOSE} \\
            }
        }
    }
    & & \\ \\

    \theutterance \stepcounter{utterance}  
    & & \multicolumn{4}{p{0.6\linewidth}}{
        \cellcolor[rgb]{0.95,0.95,0.95}{
            \makecell[{{p{\linewidth}}}]{
                \texttt{\tiny{[P1$\leftrightarrow$P]}}
                \texttt{Post{-}episode Reflexion (not sent to GM):} \\
\texttt{New plan: I failed because my previous limerick did not strictly include the word "limerick" within the body of the poem itself, only in the title or context, and I may have missed a subtle punctuation error. Next time, I will first repeat the prompt exactly as requested. Then, I will compose a five{-}line limerick that explicitly contains the words "limerick" and "funny" within the lines themselves. I will rigorously scan the final output to ensure zero commas are present, replacing any potential pauses with line breaks or alternative phrasing, and verify the AABBA rhyme scheme is intact before submitting.} \\
            }
        }
    }
    & \\ \\

\end{supertabular}
}

\end{document}
